\scanpage{24}
\uline{Proposition} Considérons une face $f=\{u,v,w\}$. \deleted{Alors les trois arêtes} Soit $t_f=t$ l'ens des trois arêtes de $f$ (en correspondance $1$-$1$ avec l'ens de ses sommets), et soient $\{u_1,v_1,w_1\}$ les troisièmes sommets des \inserted{trois} faces \inserted{$\ne f$} adjacentes à ces arêtes. Alors

a) \deleted{\illegible{1}} $\rho$ induit la permutation circulaire \unsure{standart|standard} sur \unsure{tri-ens|les ens|les} \inserted{de} directions d'arêtes $\overline{\{u,v_1\}},\ \overline{\{v,w_1\}}\quad\overline{\{w,u_1\}}$ et sur l'ens des directions d'arêtes $\overline{\{u,w_1\}},\ \overline{\{w,v_1\}},\ \overline{\{v,u_1\}}$,
\margin{$u\,v\,w$\\$u\,w\,v$}
\deleted{et} \inserted{Considérons} les trièdres \inserted{dist.} associés à ces orbites de $\rho$, en corr. $1$-$1$ avec l'ens \inserted{$\{\omega,\omega'\}=\{\omega,\omega^{-1}\}$} des deux orientations (i.e des 2 ordres circulaires) sur $f$, soit $\{t_\omega,t_{\omega'}\}$

b) Les\deleted{\unsure{ens|trois}} trièdres distingués $\mathrm{tr}\,\overline{\{v,w\}}=t_u$, $\mathrm{tr}\,\overline{\{w,u\}}=t_v$, $\mathrm{tr}\,\overline{\{u,v\}}=t_w$, $t_\omega$ et $t_{\omega'}$ sont distincts.

\uline{Dém} a) est clair. \deleted{\unsure{Que|Pour}} $t_u,t_v,t_w$ soient distincts est un cas particulier de la prop. précédente, ainsi que $t_u\ne t_\omega=\mathrm{tr}\,\overline{\{w,u_1\}}$ $t_u\ne t_{\omega'}=\mathrm{tr}\,\overline{\{v,u_1\}}$ et de $=$ pour $t_v,t_w$.

\uline{Corollaire} On a une bijection canonique
\[
T\simeq f\amalg\mathrm{Or}(f)
\]
% archive sheet number: 23
